\chapter{执行成本、资源预算与测量}
\label{ch:performance}

本章描述当前 \code{crates/core/src/verify.rs}、\code{api.rs} 的计数与实现，以及仓库已保存的受控实验。它不把 D0 的 streaming tape、8-byte LiveFact、mmap 或 batch 设想当作现有性能。当前 CLI 读取 owned bytes，decoder 构建完整 owned model/certificate，再执行 verifier；只读输入、已解码对象、精确算术临时量和 proof scopes 均会占内存。

\section{公开计数的确切含义}

\code{VERIFIED} report 的 \code{statistics} 由 kernel 生成，不是作者提交的证书字段。下面列出主要计数；全部字段定义在 \code{Statistics}。

\begin{longtable}{p{0.31\textwidth}p{0.61\textwidth}}
\toprule
字段 & 当前计数边界 \\
\midrule
\endhead
\code{obligations} & 实际执行的 terminal/output checks，包括方向 Return、严格改进和 injectivity 等通过同一 obligation 入口的检查。Split 的两支均计数；不含 lift premises、local Cases Witness、point checks 或 fiber 内部检查。\\
\code{linear_terms} & Lin 展开所访问 premise expressions 的非零 coefficient 项数，不是 Lin 数或 scalar arithmetic 次数。\\
\code{substitution_terms} & substitute 展开每个被引用 map expression 时访问的非零项总数；不同 rows 的重复访问重复计数。\\
\code{branches} & ordinary Split 与 local Cases 的 branch points；一个二支 split 计一个 point，而非两个 leaf。\\
\code{peak_live_facts} & 同一时刻当前 scope 可见的 logical local-register bindings 的峰值；不是累计定义数，也不是物理存储或 RSS。\\
\code{work} & 被明确 charge 的操作总量，包括 step、expression/fact/map 项及常数、ancestor lookup probes、构造器和相关验证工作。是资源政策单位，不是 CPU instructions、运行时间或严格 bit-complexity 上界。\\
\code{points}, \code{point_coefficient_terms} & 完整 point checks 及 point/observation 精确 evaluation 项数；不含 fiber 反证中的 Lin/代换项。Point 的 bounds/rows 检查不进入方向 Return 的 obligations。\\
\code{plotkin_pairs} & Plotkin witness 中已检查的 coordinate pairs；还存在 labels、membership 与 premise 扫描成本。\\
\code{flow_*}, \code{interval_*}, \code{menu_*}, \code{signed_*}, \code{subflow_*} & 各类 lift 次数及其 generated variables、rows、terms。记录完整 generated extension，不代表存在性引理已被另一个 universal checker 证明。\\
\code{bijections}, \code{injectivity_*} & V7 paired-domain bijection 与其完整 injectivity domain/checks。Legacy \code{inverse} 的成功不自动增加此 bijections 计数。\\
\code{objective_level_*}, \code{halfspace_*}, \code{improvements} & 条件域构造次数、生成行和项，以及严格改进 witnesses；V10 的新增 zero counters 会从 JSON 中省略。\\
\bottomrule
\end{longtable}

Context 的 \code{charge} 做 checked addition，超过有效 work ceiling 就返回 \code{RESOURCE_LIMIT}。例如一次 static/local fact 读取还计入 expression 的 terms+1；local 引用另外计 ancestor 查找 probes。一次 map substitution 按引用表达式的 terms+1 charge。换一种内部计数政策可能改变 work，必须在实验中说明；同一 work 不能保证同一机器上耗时相同。

第\ref{ch:examples}章的分段证书只有一组三项 Forward 数学要求，但 verifier 在两个 Return 各检查一次，所以 Forward 的 obligations 是六。完整 fiber refutation 使用同一 verifier 的 fresh-scope 调用，但 \code{point.rs} 只把其 work 加回共享预算，不把内部 branches、linear\_terms、substitution\_terms、peak\_live\_facts、plotkin\_pairs 或 obligations 汇入顶层 Statistics；分析这类 proof 时不能从顶层零计数断言内部没有该工作。无限精度 rationals 的位长不同，也会使同样项数的 Lin 成本很不同。不能把任何一个单项计数当作完整成本。

\section{共享预算与可复核的资源结果}

当前默认 \code{Limits} 为 work $10^9$、live facts $10^6$、depth 64，root depth 为 1。一个 request 的全部方向、条件域、lift、injectivity、points/fiber refutations 和 improvements 共用 work；嵌套 scope 或新方向不会重置。CLI 可配置 work 与 live facts，depth 不提供 CLI override。

\begin{lstlisting}[caption={预算是实际 verifier 参数}]
target/release/rpx verify --work 1000000000 --live-facts 1000000 cases/test58/source.mcm cases/test58/target.mcm cases/test58/proof.rpx cases/test58/request.json
\end{lstlisting}

有限预算是部署政策，不是模型数学定义。\code{INVALID}、\code{RESOURCE_LIMIT} 和 \code{IO_ERROR} 都不建立 reformulation 不成立。第\ref{ch:examples}章已在同一正例上实际观察 \code{--work 0} 的 RESOURCE\_LIMIT；test36 的同一完整产物在 $10^9$ work 政策下成功，旧 $10^8$ ceiling 下耗尽预算。这不是数学关系改变，也不能用失败运行时间作 speedup 的分母。

还必须分别记录其他层次的边界：CLI 每文件 owned read 为 512 MiB，RPA/request 文本为 256 MiB，serde JSON nesting capacity 为 128，binary nesting 为 64；模型 dimensions、counted arrays 和 rational magnitude 有各自上限。proof depth 与 JSON wrappers 的嵌套数不同。当前精确算术中间 numerator/denominator 的 bit budget 为 65,536。提高 work 不会绕过这些独立边界。

\section{代换、精确累加与表示成本}

对 output row $r$ 和完整 map $y_j=E_j(x)$，直接代换展开
\[
A_rE(x)-u_r=c_r+\sum_{j\in\operatorname{supp}(A_r)}a_{rj}E_j(x),
\qquad
W_r=\sum_{j\in\operatorname{supp}(A_r)}\operatorname{nnz}(E_j).
\]
输入项可能重合、相消，输出可能变稠密；$W_r$ 不只由 output row 的 nnz 决定。跨 row 重复的代换仍付出访问成本。当前 \code{Accumulator} 用 hash map 暂存精确 coefficient，再排序生成 canonical sparse expression；hash 只选择容器 bucket，精确 key/rational equality 决定接受。没有通用 affine-expression DAG、共享 mapping DAG 或证书中间结果缓存的格式承诺。

\code{BigRational} 保持精确数学语义。现有 denominator-one、乘 $\pm1$ 和零加值快路径避免部分通用算术工作，同时保留所需 intermediate/result checks；非 unit rational 仍使用 arbitrary-precision 运算。它不等于 D0 建议的机器小整数与大整数自动提升 ABI，也不把 coefficient 位长或 gcd 成本计成 CPU 指令。

展开长 prefix maps 可有二次物理大小；scalar injectivity proofs 也可能很大。二进制减少 JSON 文本开销，并未保证完整 proof resident memory 与 peak live facts 成正比。一个解释性的成本分解是
\[
T=T_{\rm read}+T_{\rm decode}+T_{\rm execute}+T_{\rm process},
\qquad
T_{\rm execute}=T_{\rm algebra}+T_{\rm scopes}+T_{\rm builders}+T_{\rm checks}.
\]
它不是已证明的统一渐近上界；精确算术位长与保存的 physical representation 需要额外描述。

\section{Shared scopes 的实现与内存边界}

当前 \code{Facts} 借用不可变 parent，只在 child 保存 local definitions 与 deleted-ID masks。其可见集合为
\[
L\cup(P\setminus D),
\]
其中 local bindings 优先；drop/reuse/drop 不会复活祖先旧值。一个 child 的 DROP 不改变 sibling 或 parent。Cases 的两支独立检查，仅把共同结论插入 parent；ordinary Split 没有后续 continuation，left 检查后 right 可接管未改变的 parent 状态。这消除了每个 branch 复制全部 inherited facts 的成本。

祖先 lookup 的 probes 被计费，深度仍有界；local insertion 和 deletion-mask growth 使用 fallible reservation。\code{peak_live_facts} 只计当前可见 bindings。被 child mask 隐藏的 immutable ancestor payload 仍在内存中，masks 也占空间，所以 logical live limit 不能直接转换成固定 bytes/live-fact 或 peak RSS 上限。decoder 的 owned objects、arithmetic allocations 及部分依赖库 allocation failure 也是当前内存保证的边界。

独立 copied-BTreeMap state oracle 的 127,072 次操作、depth-64 checks 与 negative scope tests 见 \code{audit/SHARED_SCOPE_REVIEW.md}；生产集成另见 \code{audit/scope_v6_merge.json}。状态正确性审阅和 latency 测量分别保存，不能互相代替。

\section{时钟字段与端到端成本}

当前 Rust façade 和 CLI 使用同一报告函数，公开 microseconds 字段：
\begin{longtable}{p{0.28\textwidth}p{0.64\textwidth}}
\toprule
字段 & 含义 \\
\midrule
\endhead
\code{read_us} & 读取 source、target、RPX 和 request；filesystem-free \code{verify_bytes} 不返回此字段。\\
\code{model_decode_us} & 两份完整 MCM 的 decoding/validation 阶段。\\
\code{proof_decode_us} & RPX decoding 加 caller request text parsing。\\
\code{execute_us} & 完整 \code{verify_request_with_diagnostics} 调用，包含全部请求验证。\\
\code{total_us} & kernel API 本次调用的内部总时间，包括报告前的操作；不含 process startup/shutdown。\\
process wall time & 外部计时，另含进程启动和退出。不是另一个 kernel 公开计数字段。\\
\bottomrule
\end{longtable}

MPS/LP import、作者生成 RPA、assembler 和任何外部 oracle 必须单列。\code{CANONICALIZED}/\code{ASSEMBLED} 的 \code{total_us} 不能当作 verification time。未重测的历史 preparation 应标为 archived/unmeasured；raw 数据中的空 preparation array 或 sentinel 0 不能解释为零成本。

可靠的对比先固定完整数学 request、模型/proof bytes、source/binary/lockfile identities、硬件、编译选项和 budget，再记录 cold/warm 文件状态、重复次数、顺序与 isolation。每个样本都核对退出码和 VERIFIED，保存失败和 outliers；timing 与 RSS 另起 processes。至少报告 min/median/max 或完整分布，并同时保存计数、模型 dimensions、artifact size 与阶段时间。

\section{已有实验：明确保留适用范围}

\subsection{共享状态与精确算术}

\code{bench/SHARED_SCOPE_REPORT.md} 固定了 frozen V5 的 before/after sources，数学规则和 certificate bytes 相同，仅改变 scope 表示及相应 work accounting。三个 copying-heavy workloads 的 execution medians 如下，单位 ms：
\begin{center}
\begin{tabular}{lrrrr}
\toprule
active facts / Cases & before & after & work before & work after \\
\midrule
256 / 256 & 3.439 & 0.425 & 267,535 & 8,720 \\
1,024 / 1,024 & 51.413 & 1.605 & 4,215,823 & 34,832 \\
4,096 / 4,096 & 827.142 & 6.517 & 67,194,895 & 139,280 \\
\bottomrule
\end{tabular}
\end{center}

这支持消除所选二次复制成本，不支持普遍 127 倍加速。完整 broader matrix 保存了 nonbranching regressions：例如 real65 execution median 从 3690.334 增到 3793.140 ms；logical counters 中除了 work 以外的数学项、全部 physical artifacts 一致。test18 后来的成功是原字节的预算内 replay，其原失败没有成功性能分母。

\code{bench/RATIONAL_SHORTCUTS_REPORT.md} 对 49 个 paired workloads、1,176 个成功 verifier processes 观察同 bytes、同完整计数下的算术快路径。所有 execution medians 降低，但 startup 主导的 process latency 不全部降低；例如 injectivity-64 从 2.186 增到 2.391 ms。RSS median 变化在 3\% 内。报告描述其固定 prototype binary，不能把后来 resource/API 修订后的可执行文件自动等同于被测候选。生产合并边界与额外检查记录在 \code{audit/V7_ARITHMETIC_MERGE.md}。

\subsection{V10 旧任务 regression 与顺序敏感性}

\code{bench/V10_REGRESSION_REPORT.md} 保留 31 个 paired 旧任务、712 个成功 verification/RSS processes，相同 artifacts、request 和完整 non-timing results。candidate 在 19 个任务的 execution median、23 个任务的 full-process median 上增加。test65 execution 为 1792.479 到 1864.392 ms，约 $+4.0\%$；rational-1000 为 0.792 到 0.972 ms，增加 0.180 ms。没有通用加速或无 regression 声明。

预先声明的一次 alternating-order follow-up 保存另外 72 个 timed calls 和八个 excluded warmups，四个任务的 paired median difference 分别为 rational-1000 $+0.003$ ms、test65 $+82.552$ ms、test71 $-0.038$ ms、test79 $-0.005$ ms。test65 的正成本保留为交付限制，其他三项原有较大差异未在本轮持续。分布和 within-pair order effects 不支持把一个精确百分比归因于代码。原矩阵未被替换，也没有新的 RSS 声明；原始报告见 \code{bench/results/v10-order-sensitivity/ORDER_SENSITIVITY_REPORT.md}。

\subsection{新能力绝对 baseline}

\code{bench/NEXT_CAPABILITY_DRAFT_BASELINE.md} 固定 22 个任务，264 个成功 verification/RSS processes，并单列 63 次 synthetic author/assembly preparations。使用最终 cold checkpoint 的同一 binary，task-owned build/preparation/acceptance 工作在采样期间 idle；warm files、fresh process，OS activity 未控制。新任务不与无法证明相同命题的旧 binary 计算 speed ratio。

两个已有报告的 median 样本如下；完整 min/median/max、decoder 时间和身份都在原报告：
\begin{center}
\begin{tabular}{lrrr}
\toprule
任务 & execute ms & process ms & peak RSS MiB \\
\midrule
完整 actual test58 & 18.369 & 28.749 & 31.891 \\
subflow 1,024 nodes / 4 groups & 56.732 & 87.707 & 69.141 \\
\bottomrule
\end{tabular}
\end{center}

第二项的完整 Python preparation median 为 9.154 s，包含简单 $O(nm)$ author construction；不能把这个成本藏在 constant-time theorem 的叙述中。actual test58 的历史 author preparation 未重测，表中没有零成本含义。Corpus 的总 functional replay 耗时和本手册单次示例 replay 都不是该受控矩阵的替代。

\section{复现与下一次优化的登记}

仓库已有 fixed matrix drivers；新的运行须用未使用的 label 保存 raw evidence。下面的命令可建立当前 binary 的新绝对测量，不会复写历史样本：
\begin{lstlisting}[caption={当前完整初始与 extended matrices}]
python3 tools/benchmark.py --kernel target/release/rpx --source-root . --label manual-current-initial
python3 tools/benchmark_extended.py --kernel target/release/rpx --source-root . --label manual-current-extended
\end{lstlisting}

正式执行前应安排独立采样窗口并登记机器与其它任务状态；这些命令不自动保证 isolation。固定 V10 final-driver 与能力 config 还验证专门 binary/manifest hashes，因此重建当前 binary 不会被默认为同一历史样本。已有 driver 拒绝覆盖 raw evidence，应换新 label 或单独建立新固定计划。

共享 DAG、dedicated batch、streaming decoding、模型缓存和独立 solution reconstruction 是可能的后续工作，均非当前报告中的既有收益。新增 batch 必须另定义可观察语义、operand 检查、scope/obligation 与失败条件，不能用名称加入隐藏 theorem。优化接受路径继续做精确验证；浮点 residual、probabilistic hash 和外部 solver success 不构成接受条件。

任何对其它 proof format 的速度比较，都须提供相同数学请求和推导、相同线程政策及 preparation/verification 边界。当前报告没有这种受控 cross-format 实验，因而不宣称相对 VIPR 或求解器的性能优势。

历史 source snapshots 与部分 prototype 完整副本已在交付后的清理中删除，
原 manifests、receipts、logs 与 diffs 保留；具体清理边界见
\code{audit/POST_DELIVERY_CHANGES.md}。依赖这些副本的历史重建必须先恢复原版本数据。
本章的新测量入口使用当前源码，不能把重新编译的当前 binary 当作原实验 binary。
